\section{Results}
\label{sec:results}

\subsection{Headline: the bundled method as specified}

Table~\ref{tab:main} is generated directly by \texttt{analysis/aggregate.py} from
\texttt{results/summary.csv}; no number in it is hand-typed.

\begin{table}[htbp]
\centering
\begin{tabular}{lrrrrrr}
\toprule
 & \multicolumn{3}{c}{Accuracy} & \multicolumn{3}{c}{ECE} \\
\cmidrule(lr){2-4}\cmidrule(lr){5-7}
 & \texttt{clipzs} & \texttt{tda} & \texttt{cal\_tda} & \texttt{clipzs} & \texttt{tda} & \texttt{cal\_tda} \\
\midrule
Aircraft & 23.85 & 24.27 & 23.97 & 5.44 & 15.51 & 8.65 \\
DTD & 44.39 & 47.04 & 45.63 & 8.27 & 17.41 & 17.72 \\
EuroSAT & 42.06 & 53.17 & 48.88 & 7.16 & 14.38 & 17.53 \\
Flowers102 & 67.32 & 69.91 & 68.86 & 2.63 & 6.39 & 32.73 \\
Pets & 88.20 & 88.77 & 88.74 & 4.37 & 2.52 & 43.98 \\
\bottomrule
\end{tabular}
\caption{Top-1 accuracy and ECE (\%, 20 bins) on the five test splits.}
\label{tab:main}
\end{table}


Averaged over the five datasets, zero-shot CLIP reaches 53.16\% accuracy at 5.57\% ECE, TDA
reaches 56.63\% at 11.24\% ECE, and \texttt{cal\_tda} with all three contributions enabled --
the method exactly as specified -- reaches 55.22\% at 24.12\% ECE. \textbf{This is a clear
failure of the bundled method.} Accuracy is lower than TDA's on every single dataset (from 0.03
points on Pets up to 4.29 points on EuroSAT), and ECE is worse than TDA's on four of the five
datasets, dramatically so on Flowers102 (32.73\% vs.\ 6.39\%) and Pets (43.98\% vs.\ 2.52\%).
Aircraft is the sole dataset where the bundle's ECE number looks better than TDA's (8.65\%
vs.\ 15.51\%) -- Section~\ref{sec:discussion} shows this is not real correction, only the same
underconfidence failure landing closer to correct by accident on this one low-accuracy dataset.
We do
not reframe this as a partial success: on the spec's own success criterion (accuracy $\geq$ TDA
\emph{and} ECE $<$ TDA, on at least 3 of 5 datasets), the bundled method meets the bar on
\textbf{0 of 5} datasets.

Figure~\ref{fig:scatter} plots the same comparison directly: our success criterion is a point as far right as
TDA (equal accuracy) but as far down as zero-shot CLIP (honest confidence). The bundled method's
green star sits left of TDA's orange square on every dataset (lower accuracy) and above it on
four of the five (higher ECE) -- Aircraft is the one dataset where the star sits below the
square, matching the exception already noted above.

\begin{figure}[htbp]
    \centering
    \includegraphics[width=0.75\textwidth]{accuracy_vs_ece.png}
    \caption{Accuracy vs.\ ECE, all five datasets. The purple diamond (LOO temperature only,
    Section~\ref{sec:ablation}) is the configuration this report ultimately recommends, plotted
    alongside the three headline methods. Marker shape and colour identify the method (see legend);
    each dataset's four points are joined by a line in that dataset's own colour, labelled at the
    line's right-hand end. Lines run left to right by accuracy; they mark grouping only, not a
    progression.}
    \label{fig:scatter}
\end{figure}

\subsection{Reliability diagrams}

Figure~\ref{fig:reliability} plots, per dataset, the confidence-vs-accuracy curve for each of
the three headline methods plus the ablation's \texttt{temp}-only variant (a curve on the
dashed identity line is perfectly calibrated). The
grid makes the aggregate ECE numbers in Table~\ref{tab:main} legible per bin, and the direction
of the failure is visible directly: on every dataset, \texttt{cal\_tda}'s curve sits
\emph{above} the identity line -- accuracy exceeds stated confidence, not the reverse. In
the five \texttt{cal\_tda} reliability records the highest occupied confidence bin tops out at
0.70 on EuroSAT and at 0.60 or below on the other four, so there is no high-confidence region
to be confidently wrong in; the method is instead severely \emph{underconfident}, most visibly on
Pets, where 88.74\% accuracy is reported at only 44.77\% mean confidence. Section~\ref{sec:discussion}
explains the mechanism, including why Aircraft -- the one dataset where the bundle's ECE
looked like an improvement over TDA in Table~\ref{tab:main} -- is not a case of the method
correcting itself, but the smallest instance of the same underconfidence effect.

\begin{figure}[htbp]
    \centering
    % Three panels per row. Two-per-row was tried and reverted: five square panels
    % in three rows exceed \textheight, and LaTeX responds by pushing the caption
    % and the last subcaption clean off the bottom of the paper ("Float too large
    % for page by 128pt"). On-page legibility is fixed in make_figures.py instead,
    % by shrinking the matplotlib canvas rather than enlarging the float. Slot six
    % holds the legend the five panels share.
    \begin{subfigure}[t]{0.32\textwidth}
        \includegraphics[width=\linewidth]{reliability_dtd.png}
        \caption{DTD}
    \end{subfigure}
    \begin{subfigure}[t]{0.32\textwidth}
        \includegraphics[width=\linewidth]{reliability_flower102.png}
        \caption{Flowers102}
    \end{subfigure}
    \begin{subfigure}[t]{0.32\textwidth}
        \includegraphics[width=\linewidth]{reliability_pets.png}
        \caption{Pets}
    \end{subfigure}\\[1ex]
    \begin{subfigure}[t]{0.32\textwidth}
        \includegraphics[width=\linewidth]{reliability_eurosat.png}
        \caption{EuroSAT}
    \end{subfigure}
    \begin{subfigure}[t]{0.32\textwidth}
        \includegraphics[width=\linewidth]{reliability_aircraft.png}
        \caption{Aircraft}
    \end{subfigure}
    % A minipage, not a subfigure: the legend is not a sixth panel and must not be
    % lettered (f) or counted in the figure's subcaptions.
    %
    % [b], not [t], and an explicit height. Given a height, minipage's [t] puts the
    % box *top* on the row's baseline -- which is the bottom edge of the panels
    % beside it, so the legend landed a full panel below the row. [b] puts the box
    % bottom there instead, so a box exactly one panel high (0.32\textwidth: the
    % panels are square) spans the row, and inner-pos c centres the legend in it.
    \begin{minipage}[b][0.32\textwidth][c]{0.32\textwidth}
        \centering
        \includegraphics[width=\linewidth]{reliability_legend.png}
    \end{minipage}
    \caption{Reliability diagrams, all five datasets: zero-shot CLIP, TDA, \texttt{cal\_tda}
    (all three contributions), and the ablation's \texttt{temp}-only variant. All five panels
    share the legend at bottom right; a curve on the dashed identity line is perfectly calibrated.}
    \label{fig:reliability}
\end{figure}

\subsection{The ablation}
\label{sec:ablation}

Table~\ref{tab:ablation} (generated by \texttt{analysis/make\_ablation\_table.py} from a
40-run sweep -- all eight on/off combinations of the three contributions, on all five datasets)
isolates which of the three changes are responsible for the failure above.

\begin{table}[htbp]
\centering
\begin{tabular}{lrr}
\toprule
 & Accuracy & ECE \\
\midrule
none: TDA-equivalent (all flags off) & 56.60 & 11.28 \\
margin: + margin admission & 53.67 & 18.06 \\
fusion: + probabilistic fusion & 55.60 & 24.54 \\
temp: + LOO temperature & 56.60 & 8.54 \\
margin\_fusion: + margin + fusion & 55.22 & 23.94 \\
margin\_temp: + margin + temperature & 53.67 & 19.39 \\
fusion\_temp: + fusion + temperature & 55.60 & 24.78 \\
all: all three (ours) & 55.22 & 24.12 \\
\bottomrule
\end{tabular}
\caption{Ablation: mean accuracy and ECE (\%) across the five datasets.}
\label{tab:ablation}
\end{table}


\textbf{Contribution 3 (LOO temperature scaling) alone succeeds.} The ablation's own \texttt{none} baseline is
\texttt{cal\_tda} with all three flags off, run inside the same 40-run sweep; it is the correct
same-code-path comparison, since \texttt{temp} and \texttt{none} differ by exactly one config
flag. Against it, \texttt{temp} matches accuracy \emph{exactly}, to 0.00 points, on all five
datasets -- a structural guarantee, since temperature scaling is a monotonic rescaling of the
logits and cannot change an argmax. It cuts mean ECE from 11.28\% to 8.54\%, the best mean ECE
of all eight ablation variants, including \texttt{all} (24.12\%). The per-dataset breakdown is in Table~\ref{tab:temp-per-dataset}:

\begin{center}
\begin{tabular}{lrrr}
\toprule
Dataset & \texttt{none} ECE & \texttt{temp} ECE & Change \\
\midrule
Aircraft   & 15.56 & 4.21  & $-11.35$ \\
DTD        & 17.41 & 4.68  & $-12.73$ \\
EuroSAT    & 14.43 & 13.43 & $-1.00$ \\
Flowers102 & 6.51  & 4.85  & $-1.66$ \\
% \textbf does not reach inside math mode, so $+13.06$ printed at roman weight
% next to a bold "(worse)" in the same cell. \mathbf carries the weight in.
Pets       & 2.48  & \textbf{15.54} & \textbf{$\mathbf{+13.06}$ (worse)} \\
\bottomrule
\end{tabular}
\captionof{table}{Per-dataset ECE (\%) for the ablation's \texttt{none} and \texttt{temp} variants. Pets is the one reversal.}
\label{tab:temp-per-dataset}
\end{center}

This is not a uniform win: on Pets, where the \texttt{none} baseline was already very well
calibrated (2.48\% ECE),
LOO temperature scaling makes calibration substantially \emph{worse}. We discuss the mechanism
in Section~\ref{sec:discussion} rather than average it away.

\textbf{Two counts, and why they differ.} Against the ablation's \texttt{none} baseline,
\texttt{temp} meets the spec's success criterion on \textbf{4 of 5} datasets (all but Pets).
Against the separately run headline \texttt{tda} rows in Table~\ref{tab:main}, the count is
lower -- \textbf{2 of 5} -- because \texttt{none} and \texttt{tda} are two different process
invocations of what should be identical logic, and in practice differ by 0--4 flipped
predictions per dataset out of 1{,}692--8{,}100 samples. The cause is not a logic difference:
upstream \texttt{tda} accumulates its final logits in fp16, while \texttt{cal\_tda} (even
with every flag off) upcasts to fp32 before its final step, and at CLIP's logit magnitudes that
rounding difference is worth roughly one unit in the last place -- enough, on
borderline-margin predictions, to flip an argmax once every few hundred to a few thousand
samples. Both counts are correct; they
answer different questions. The 4-of-5 count answers whether temperature scaling, holding
everything else fixed, helps calibration without moving accuracy -- yes, on 4 of 5. The
2-of-5 count answers whether a \texttt{temp}-only run, as a real separate process, beats the
specific \texttt{tda} numbers quoted in Table~\ref{tab:main} -- only on 2 of 5, because of run-to-run floating-point drift that
is two orders of magnitude smaller than any effect discussed above.

\textbf{Contributions 1 and 2 are the ones actually responsible for the bundle's failure.}
Margin admission alone \emph{lowers} mean accuracy (56.60\% $\rightarrow$ 53.67\%) and raises
mean ECE (11.28\% $\rightarrow$ 18.06\%): its selectivity throws away cache signal the
unfiltered cache was using correctly on some datasets. Probabilistic fusion alone is the worst
single component: it raises mean ECE to 24.54\%, driven by blowups on Flowers102 (32.46\%) and
Pets (43.77\%) specifically, though it happens to help Aircraft (15.56\% $\rightarrow$
9.16\%) -- the reason ``all three'' still beats plain TDA's ECE on that one dataset. Neither
contribution helps consistently even in isolation; both actively hurt the mean.

Cache admission rates, which explain part of why margin admission behaves so differently across
datasets, are reported in Section~\ref{sec:admission-rates}.
